\subsection{Задание}

Реализовать информационную систему, которая позволяет взаимодействовать с объектами класса \texttt{Movie}, описание которого приведено ниже:

\begin{Verbatim}
public class Movie {
    private Long id; //Поле не может быть null, Значение поля должно быть больше 0, Значение этого поля должно быть уникальным, Значение этого поля должно генерироваться автоматически
    private String name; //Поле не может быть null, Строка не может быть пустой
    private Coordinates coordinates; //Поле не может быть null
    private java.time.LocalDate creationDate; //Поле не может быть null, Значение этого поля должно генерироваться автоматически
    private int oscarsCount; //Значение поля должно быть больше 0
    private long budget; //Значение поля должно быть больше 0
    private Double totalBoxOffice; //Поле может быть null, Значение поля должно быть больше 0
    private MpaaRating mpaaRating; //Поле не может быть null
    private Person director; //Поле может быть null
    private Person screenwriter;
    private Person operator; //Поле может быть null
    private Integer length; //Поле не может быть null, Значение поля должно быть больше 0
    private Long goldenPalmCount; //Значение поля должно быть больше 0, Поле может быть null
    private int usaBoxOffice; //Значение поля должно быть больше 0
    private String tagline; //Поле не может быть null
    private MovieGenre genre; //Поле может быть null
}
public class Coordinates {
    private Integer x; //Максимальное значение поля: 826, Поле не может быть null
    private Integer y; //Поле не может быть null
}
public class Person {
    private String name; //Поле не может быть null, Строка не может быть пустой
    private Color eyeColor; //Поле может быть null
    private Color hairColor; //Поле может быть null
    private Location location; //Поле может быть null
    private java.time.ZonedDateTime birthday; //Поле не может быть null
    private Float height; //Поле может быть null, Значение поля должно быть больше 0
    private Float weight; //Поле может быть null, Значение поля должно быть больше 0
    private String passportID; //Значение этого поля должно быть уникальным, Поле не может быть null
}
public class Location {
    private Long x; //Поле не может быть null
    private Double y; //Поле не может быть null
    private String name; //Поле не может быть null
}
public enum MpaaRating {
    G,
    PG,
    PG_13,
    R,
    NC_17;
}
public enum MovieGenre {
    WESTERN,
    COMEDY,
    TRAGEDY,
    HORROR;
}
public enum Color {
    GREEN,
    BLACK,
    ORANGE,
    WHITE,
    BROWN;
}
public enum Country {
    GERMANY,
    VATICAN,
    ITALY,
    NORTH_KOREA;
}
\end{Verbatim}

\subsection{Требования к системе}

Разработанная система должна удовлетворять следующим требованиям:

\begin{itemize}
\item Основное назначение информационной системы - управление объектами, созданными на основе заданного в варианте класса.
\item Необходимо, чтобы с помощью системы можно было выполнить следующие операции с объектами: создание нового объекта, получение информации об объекте по ИД, обновление объекта (модификация его атрибутов), удаление объекта. Операции должны осуществляться в отдельных окнах (интерфейсах) приложения. При получении информации об объекте класса должна также выводиться информация о связанных с ним объектах.
\item При создании объекта класса необходимо дать пользователю возможность связать новый объект с объектами вспомогательных классов, которые могут быть связаны с созданным объектом и уже есть в системе.
\item Выполнение операций по управлению объектами должно осуществляться на серверной части (не на клиенте), изменения должны синхронизироваться с базой данных.
\item На главном экране системы должен выводиться список текущих объетов в виде таблицы (каждый атрибут объекта - отдельная колонка в таблице). При отображении таблицы должна использоваться пагинация (если все объекты не помещаются на одном экране).
\item Нужно обеспечить возможность фильтровать/сортировать строки таблицы, которые показывают объекты (по значениям любой из строковых колонок). Фильтрация элементов должна производиться только по полному совпадению.
\item Переход к обновлению (модификации) объекта должен быть возможен из таблицы с общим списком объектов и из области с визуализацией объекта (при ее реализации).
\item При добавлении/удалении/изменении объекта, он должен автоматически появиться/исчезнуть/измениться в интерфейсах у других пользователей, авторизованных в системе.
\item Если при удалении объекта с ним связан другой объект, связанные объекты должны удаляться.
\item Пользователи должны иметь возможность просмотра всех объектов. Для модификации объекта должно открываться отдельное диалоговое окно. При вводе некорректных значений в поля объекта должны появляться информативные сообщения о соответствующих ошибках.
\end{itemize}

\subsection{Специальные операции}

В системе должен быть реализован отдельный пользовательский интерфейс для выполнения специальных операций над объектами:

\begin{itemize}
\item Удалить все объекты, значение поля usaBoxOffice которого эквивалентно заданному.
\item Рассчитать среднее значение поля usaBoxOffice для всех объектов.
\item Вернуть массив объектов, значение поля tagline которых начинается с заданной подстроки.
\item Равномерно перераспределить все "Оскары", полученные фильмами одного жанра, между фильмами другого жанра.
\item Дополнительно наградить все фильмы с продолжительностью больше заданной указанным числом "Оскаров".
\end{itemize}

Представленные операции должны быть реализованы в рамках компонентов бизнес-логики приложения без прямого использования функций и процедур БД.

\subsection{Особенности хранения объектов}

\begin{itemize}
\item Организовать хранение данных об объектах в реляционной СУБД (PostgreSQL). Каждый объект, с которым работает ИС, должен быть сохранен в базе данных.
\item Все требования к полям класса (указанные в виде комментариев к описанию классов) должны быть выполнены на уровне ORM и БД.
\item Для генерации поля id использовать средства базы данных.
\item Для подключения к БД на кафедральном сервере использовать хост pg, имя базы данных - studs, имя пользователя/пароль совпадают с таковыми для подключения к серверу.
\end{itemize}

\subsection{Организация взаимодействия с пользователем}

При создании системы нужно учитывать следующие особенности организации взаимодействия с пользователем:

\begin{itemize}
\item Система должна реагировать на некорректный пользовательский ввод, ограничивая ввод недопустимых значений и информируя пользователей о причине ошибки.
\item Переходы между различными логически обособленными частями системы должны осуществляться с помощью меню.
\item При добавлении/удалении/изменении объекта, он должен автоматически появиться/исчезнуть/измениться на области у всех других клиентов.
\end{itemize}

\subsection{Требования к разработке}

\begin{itemize}
\item В качестве основы для реализации ИС необходимо использовать Java EE, Managed Beans.
\item Для создания уровня хранения необходимо использовать EclipseLink.
\item Разные уровни приложения должны быть отделены друг от друга, разные логические части ИС должны находиться в отдельных компонентах.
\end{itemize}
